\lecture{1}{Sun 24 Apr 2022 03:55}{Proof of te problems in syllabus}

\begin{lemma}[Equality test not computable]
	The equality relation  $=$ on $\R$ is not decidable.
\end{lemma}
\begin{proof}
	Suppose there exists a representation $\delta:\subseteq \Sigma^\omega \to \R$ where the equality relation is $\delta$-decidable. Consider a real number $x$ and some $p$ such that $\delta(p) = x$. On input $\langle p, p\rangle$, the Turing machine will conclude that the inputs are equal, after a finite number of steps. Let $w$ be the prefix of $p$ that the machine has read before it stops. We know $w \in \Sigma^*$, i.e. is a finite string. Now for any other real number $x'$ with  $\delta(p') = x'$, the machine will conclude that  $x = x'$ whenever $p' \in w \Sigma^\omega$, i.e. $p'$ has $w$ as a prefix. 

	Let $\sim$ be the equivalence relation on $\Sigma^\omega$ generated by the machine, where $p \sim p' \iff $ the machine concludes after a finite number of steps that $\delta(p) = \delta(p')$. The equivalence classes $\Sigma^\omega / \sim$ are determined by finite strings in $\Sigma^*$, however $\Sigma^*$ is countable, so $\Sigma^\omega / \sim$ is countable. Since $\delta$ is onto, we conclude that  $\R$ is countable, a contradiction.
\end{proof}

\begin{lemma}
	Comparison with 0 is not computable for reals.
\end{lemma}
\begin{proof}
	$x = y \iff x - y = 0$, so this is equivalent to saying that equality test is decidable, which is false by the last lemma.
\end{proof}
\begin{lemma}
	Comparison with 0 is not computable for constructive reals.
\end{lemma}
\begin{proof}
	This one I do not know how to prove for general $\delta$. Take $\delta$ to be the Cauchy representation. Take a sequence of rational numbers that converge rapidly to 0, i.e. $|x_i|<2^{-i}$. Denote the terms by $x_1, x_2, \ldots$. This sequence generates a Cauchy representation of 0. Suppose the comparison is computable, then the Turing machine will stop after reading finite $n$ terms from the Cauchy sequence, and conclude that the number it represents is exactly 0. 

	Now any other constructive real $y \in \bigcap_{i=1}^n (x_i-2^{-i}, x_i + 2^{-i})$ can be represented by the same $x_1, \ldots, x_n$, but not necessarily equal to 0. We know such a $y$ exists because  $\R_c$ is dense in  $\R$.
\end{proof}



\begin{theorem}[problems in syllabus]
	Consider a linear equation $Ax+By+C=0$ where $A,B,C \in \R_c$. There is no algorithm that tells if the equation has a root, if it has exactly one root, or on which endpoint of a closed interval of $x$ it attains maximum for $y$.
\end{theorem}
\begin{proof}
	The linear equation has no root iff  $A = B = 0$ and  $C \neq 0$. It has infinitely many roots iff $A = B = C = 0$. It has exactly one root iff either  $A \neq 0$ or $B \neq 0$. The RHS of the equivalences above are simply comparison with 0, which are not computable by the last lemma.

	If we restrict the domain to $\{\left(x,y  \right) \in \R_c^2 : a\le x\le b \}$, it can be easily shown that the endpoint on which $y$ attains maximum only depends on sign of the gradient $\frac{dy}{dx} = - \frac{A}{B}$. Under the Cauchy representation, the arithmetic operations are computable, hence $-\frac{A}{B}$ is computable. The problem then reduces to a comparison with zero, which is not computable by the last lemma. 
\end{proof}
